%==============================================================================
%  ICVGIP-2026 SUBMISSION -- main manuscript (official ACM/acmart ICVGIP template).
%------------------------------------------------------------------------------
%  All acmart support files are bundled in this folder (from the official
%  ICVGIP-2026 template zip), so this compiles with NO package installation:
%      pdflatex main.tex ; bibtex main ; pdflatex main.tex ; pdflatex main.tex
%  (or: latexmk -pdf main.tex).
%
%  * Keep `review=true` for the review submission (per the official template).
%    For camera-ready, remove it and fill \acmDOI / \acmISBN.
%  * Replace "Submission Id XYZW" with your ICVGIP-2026 submission id.
%
%  Scientific-honesty policy (unchanged): NO results/metrics/contribution-outcome
%  claims are written until their supporting evidence exists (gates G0--G8).
%==============================================================================

\documentclass[sigconf, review=true]{acmart}
%Do not remove the review=true option for papers submitted for review to ICVGIP2026.

\usepackage{booktabs} % formal tables (ACM style)

\setcopyright{rightsretained}

\acmConference[ICVGIP2026]{17th Indian Conference on Computer Vision, Graphics and Image Processing}{December 21-25, 2026}{Kolkata, India}
\acmYear{2026}
\copyrightyear{2026}
\acmPrice{15.00}
% \acmDOI{...}    % camera-ready only
% \acmISBN{...}   % camera-ready only

% convenience macros
\newcommand{\absrel}{\mathrm{AbsRel}}
\newcommand{\tae}{\mathrm{TAE}}

\begin{document}

\title{Occlusion-Aware Instance-Level Video Depth on a Single RGB-D Sensor:
       Adaptations and a Temporal-Consistency Extension}
%  Framing: ADAPTATION + EXTENSION (locked).

\author{Submission Id XYZW}   % <-- replace XYZW with your ICVGIP-2026 submission id
\affiliation{%
  \institution{Anonymous (submitted for review)}
  \country{}
}
\renewcommand{\shortauthors}{}

\begin{abstract}
\emph{[Abstract -- written after the contribution set and results are locked.
Intentionally left blank so no unverified claim enters the manuscript
prematurely.]}
\end{abstract}

% CCS concepts -- regenerate at camera-ready via https://dl.acm.org/ccs.cfm
\begin{CCSXML}
<ccs2012>
 <concept>
  <concept_id>10010147.10010178.10010224.10010225.10010228</concept_id>
  <concept_desc>Computing methodologies~Reconstruction</concept_desc>
  <concept_significance>500</concept_significance>
 </concept>
 <concept>
  <concept_id>10010147.10010178.10010224.10010240.10010242</concept_id>
  <concept_desc>Computing methodologies~Scene understanding</concept_desc>
  <concept_significance>500</concept_significance>
 </concept>
 <concept>
  <concept_id>10010147.10010178.10010224.10010240.10010241</concept_id>
  <concept_desc>Computing methodologies~Image segmentation</concept_desc>
  <concept_significance>300</concept_significance>
 </concept>
</ccs2012>
\end{CCSXML}
\ccsdesc[500]{Computing methodologies~Reconstruction}
\ccsdesc[500]{Computing methodologies~Scene understanding}
\ccsdesc[300]{Computing methodologies~Image segmentation}

\keywords{monocular depth estimation, metric depth, instance segmentation,
  occlusion reasoning, video depth consistency, RGB-D}

\maketitle

%==============================================================================
%  BODY
%==============================================================================

\section{Introduction}
\label{sec:intro}
Estimating per-pixel depth from images underpins many downstream tasks---scene
understanding, robotics, and augmented reality---and modern monocular estimators
produce high-quality dense depth from a single
image~\cite{ranftl2021dpt,yang2024depthanythingv2}. Scenes containing several
interacting people are, however, particularly difficult: where one person
occludes another it is ambiguous which surface a pixel belongs to, and a
per-pixel estimator has no explicit notion of \emph{which} person is where.
Reasoning at the level of \emph{instances}---predicting, for each person, both a
segmentation and a depth---and enforcing that occluders lie in front of the
people they occlude is therefore attractive for group scenes.

InstanceDepth~\cite{liang2025instancedepth} recently formalized this as
instance-level, occlusion-aware depth and introduced a three-stage architecture
for it. Its design targets a particular acquisition setting; deploying the same
idea on a \emph{single calibrated RGB-D sensor} recording everyday human
activity---the setting of this paper---changes the problem in three concrete
ways (Section~\ref{sec:problem}): supervision is truly metric but available only
on \emph{visible} surfaces; the ground-truth instance masks are \emph{modal and
mutually disjoint}, so two occluding people have almost no mask overlap; and,
because depth is predicted per frame, the depth of static structure
\emph{flickers} over time. Each of these breaks an assumption that
occlusion-aware instance-depth methods---and per-frame depth methods
generally---were built on.

We study how to make occlusion-aware instance depth work in this single-sensor
regime. We reproduce the three-stage InstanceDepth pipeline, adapt each stage to
metric, visible-only, modal-mask supervision, and add a streaming module that
suppresses temporal flicker. To support the study we assemble a single-sensor
RGB-D corpus of human-activity video and annotate it automatically with instance
masks, temporally consistent identities, and per-instance depth layers.

\smallskip\noindent\textbf{Contributions.}
\begin{itemize}
  \item \textbf{Single-sensor adaptations} of occlusion-aware instance depth:
    metric (scale-variant) depth supervision in place of scale-invariant
    objectives; bounding-box occlusion pairing suited to modal, disjoint masks;
    and visible-only supervision of the occlusion refinement. We further adopt a
    frozen-base refinement that is non-degrading by construction
    (Section~\ref{ssec:adapt}).
  \item A \textbf{streaming temporal-consistency extension}---a lightweight
    recurrent depth stabilizer trained with a temporal-gradient-matching
    objective---added to the otherwise per-frame pipeline
    (Section~\ref{ssec:temporal}).
  \item A \textbf{bounded pair-attention refinement head}, which we describe and
    study against the original relation-reasoning head
    (Sections~\ref{ssec:phase3} and~\ref{sec:ablation}).
  \item A \textbf{single-sensor RGB-D human-activity corpus} with an automatic
    instance and depth-layer annotation pipeline (Section~\ref{sec:dataset}).
\end{itemize}
We evaluate each component on this corpus and analyze where it helps and where
it does not (Sections~\ref{sec:results} and~\ref{sec:ablation}).

\section{Related Work}
\label{sec:related}

\paragraph{Monocular and metric depth estimation.}
Learning-based monocular depth estimation was popularized by Eigen et
al.~\cite{eigen2014silog}, who introduced a multi-scale network and the
scale-invariant logarithmic (SILog) loss still widely used today. Subsequent
work refined the regression objective through ordinal regression over
discretized depth bins~\cite{fu2018dorn} and the reverse-Huber (berHu)
loss~\cite{laina2016berhu}. MiDaS~\cite{ranftl2022midas} showed that mixing
heterogeneous datasets under a scale-and-shift-invariant objective yields
robust \emph{relative} depth, and DPT~\cite{ranftl2021dpt} replaced
convolutional decoders with a transformer-based dense-prediction head. More
recently, self-supervised features from DINOv2~\cite{oquab2024dinov2} and the
large-scale Depth Anything~V2 model~\cite{yang2024depthanythingv2}, which
couples a DINOv2 encoder with a DPT head, have become strong monocular-depth
backbones; we build on these pretrained encoders. Unlike the scale-invariant
objectives of MiDaS and Depth Anything, our single calibrated RGB-D sensor
provides true \emph{metric} supervision, so we retain a metric (scale-variant)
SILog objective throughout.

\paragraph{Temporal consistency in video depth.}
Per-frame depth models flicker when applied to video. Video Depth
Anything~\cite{chen2025vda} addresses this for long videos with a
temporal-gradient-matching (TGM) objective that penalizes only frame-to-frame
changes absent from the ground truth, thereby avoiding the moving-object bias
of optical-flow warping losses~\cite{teed2020raft}. Recurrent operators such as
the convolutional GRU~\cite{ballas2016convgru} provide streaming,
constant-memory temporal aggregation. We adopt a lightweight streaming ConvGRU
stabilizer trained with a TGM objective as an optional extension to the
otherwise per-frame pipeline we study.

\paragraph{Instance and video instance segmentation.}
Query-based set prediction, introduced for detection by
DETR~\cite{carion2020detr}, underpins Mask2Former~\cite{cheng2022mask2former},
which performs image segmentation with masked attention and bipartite
(Hungarian) matching, using point-based supervision from
PointRend~\cite{kirillov2020pointrend} for memory efficiency. Mask2Former was
extended to video instance segmentation~\cite{cheng2021m2fvis}, and
MinVIS~\cite{huang2022minvis} showed that per-frame query embeddings are
already temporally consistent enough to track instances by embedding matching,
\emph{without} video-specific training. We use a Mask2Former instance decoder
to predict per-instance masks and depth layers, and adopt MinVIS-style query
matching for training-free instance identity across frames.

\paragraph{Promptable segmentation and tracking for annotation.}
Promptable segmentation models---Segment Anything (SAM)~\cite{kirillov2023sam}
and its video extension SAM~2~\cite{ravi2024sam2}---together with decoupled
video trackers such as DEVA~\cite{cheng2023deva}, enable scalable, weakly
supervised generation of masks and temporally consistent identities. We use a
promptable concept-level video segmenter in this lineage to annotate our
dataset automatically, replacing manual labeling of instance masks and track
identities.

% (Optional under the adaptation+extension framing) a brief mention of RGB-D
% depth benchmarks could contextualize the single-sensor setting; deferred
% until verified dataset citations are added -- no fabricated citations.

\paragraph{Instance-level occlusion-aware depth.}
Closest to our work, InstanceDepth~\cite{liang2025instancedepth} predicts
instance-level, occlusion-aware depth for groups by combining a holistic
depth-range initialization, an instance depth-layer decoder, and an
occlusion-pair relation-reasoning module, evaluated on a group-interaction
video-depth dataset. Our study adapts this method to a substantially different
acquisition setting---a \emph{single} calibrated RGB-D sensor capturing
human-activity scenes---which changes the available supervision (metric depth
on visible surfaces only, with modal and mutually disjoint instance masks) and
motivates the dataset, single-sensor adaptations, and temporal extension we
present.

\section{Problem Formulation}
\label{sec:problem}
We consider instance-level, occlusion-aware \emph{metric} depth estimation for
human-activity video acquired by a \emph{single} calibrated RGB-D sensor. A
video is a sequence of RGB frames $\{I_t\}_{t=1}^{T}$, $I_t\in\mathbb{R}^{H\times
W\times 3}$. During training each frame is paired with a metric depth map
$D^{*}_t\in\mathbb{R}^{H\times W}$ from the sensor and a validity mask
$V_t\in\{0,1\}^{H\times W}$ that is zero where the sensor returns no measurement;
per-instance modal masks $\{M^{*}_{t,k}\}_k$ with temporally consistent track
identities; and, for each instance, a ground-truth depth layer
$\mathrm{Dep}^{*}_{t,k}$ defined as the mean metric depth over that instance's
\emph{visible} pixels,
\begin{equation}
\label{eq:depthlayer}
\mathrm{Dep}^{*}_{t,k}=
\frac{\sum_{p} M^{*}_{t,k}(p)\,V_t(p)\,D^{*}_t(p)}
     {\sum_{p} M^{*}_{t,k}(p)\,V_t(p)}.
\end{equation}

Given only $\{I_t\}$ at inference, the task is to predict, for every frame $t$:
(i) a dense metric depth map $D_t$; (ii) a set of instances, each with a mask
$M_{t,k}$, a category, and an instance depth layer $\mathrm{Dep}_{t,k}$; and
(iii) occlusion-corrected instance depths $\hat{\mathrm{Dep}}_{t,k}$ that respect
the geometric ordering of overlapping instances (an occluder is nearer than the
instance it occludes). Predictions should additionally be temporally stable:
the frame-to-frame change of $D_t$ should track that of the ground truth rather
than exhibit spurious flicker.

Three properties of the single-sensor setting shape the problem and distinguish
it from the multi-view or synthetic settings in which such methods are often
developed:
\begin{itemize}
  \item \textbf{Metric but visible-only supervision.} A single sensor measures
    true metric depth, but only on surfaces it can see; there is no ground truth
    behind an occluder. Supervision for occluded regions is thus unavailable,
    and metric (rather than scale-invariant) objectives are appropriate.
  \item \textbf{Modal, mutually-disjoint instance masks.} Each pixel is assigned
    to at most one instance, so two occluding people have near-zero mask overlap
    regardless of how strongly they occlude one another; occlusion must be
    detected from bounding-box overlap rather than mask intersection.
  \item \textbf{Per-frame ambiguity and flicker.} Depth is estimated
    independently per frame, so the depth of static structure can fluctuate
    across time even when the scene does not move.
\end{itemize}
These properties motivate the single-sensor adaptations and the temporal
extension described in Section~\ref{sec:method}.

\section{Dataset and Annotation Pipeline}
\label{sec:dataset}

\subsection{Acquisition}
Our data are captured with a single calibrated ZED stereo RGB-D sensor in
indoor human-activity scenes. The corpus comprises $306$ video sequences
($50{,}323$ frames at $720\times1280$; mean $164.5$ frames per sequence),
each providing time-synchronized RGB and metric depth. All instances belong to
a single category, \emph{person}, and depth is used within the sensor's
reliable metric range $[0.01,10]\,$m. We split the corpus at the \emph{video}
level into $245$ training and $61$ test sequences ($19.9\%$ test), stratified by
capture batch and drawn with a fixed seed, so that every recording condition
appears in both splits and no frames of a sequence are shared across the split
(Table~\ref{tab:dataset}). The test split contains $10{,}400$ frames, of which
$3{,}083$ contain at least one pair of overlapping instances---the
\emph{occlusion slice} used in Section~\ref{sec:results}.

\begin{table}[t]
\centering
\caption{Dataset summary. Single-sensor RGB-D, single category
(\emph{person}), video-level train/test split.}
\label{tab:dataset}
\begin{tabular}{lccc}
\toprule
Split & Sequences & Frames & Occlusion frames \\
\midrule
Train & $245$ & $39{,}923$ & --- \\
Test  & $61$  & $10{,}400$ & $3{,}083$ \\
\midrule
Total & $306$ & $50{,}323$ & --- \\
\bottomrule
\end{tabular}
\end{table}

\subsection{Automatic annotation}
The sensor supplies RGB and metric depth but not the instance-level labels the
method requires, so we generate them automatically.
\emph{(i)~Masks and identities.} A promptable, concept-level video segmentation
model in the Segment-Anything lineage~\cite{kirillov2023sam,ravi2024sam2} is
prompted with the concept \emph{person} and propagated across each sequence,
yielding a mask and a provisional identity for every person in every frame;
this replaces the conventional two-stage mask-plus-tracker pipeline
(e.g.\ SAM followed by DEVA~\cite{cheng2023deva}) with a single pass.%
\footnote{Implementation uses a concept-prompted video segmenter; the exact
model and citation are finalized for the camera-ready.}
\emph{(ii)~Identity consolidation.} Provisional tracks are consolidated by
removing cross-prompt duplicate masks by intersection-over-union (IoU),
re-linking a track that ends to one that begins within a short temporal gap
when their boundary masks agree in IoU (a highest-IoU identity-consistency
rule), and discarding tracks shorter than a minimum length.
\emph{(iii)~Depth layers.} For each instance we compute the ground-truth depth
layer of Eq.~\eqref{eq:depthlayer}: the mean metric depth over its visible,
sensor-valid pixels.
\emph{(iv)~Single-label maps.} Masks are rasterized into one $16$-bit identity
image per frame; where instances overlap, each contested pixel is assigned to
the instance with the smaller depth layer (the nearer, occluding surface),
while the per-instance depth layers are computed from each instance's full mask
\emph{before} this resolution so that a partially occluded instance's depth is
not biased. A ground/floor mask is additionally produced by the same segmenter
for optional auxiliary supervision.

\subsection{Statistics}
On average a frame contains $2.34$ people. Per-instance depths concentrate at
close range: the large majority of instances lie within $0$--$4\,$m
($67{,}778$ within $0$--$2\,$m and $16{,}105$ within $2$--$4\,$m), with a
smaller tail toward $10\,$m ($1{,}948$; $2{,}705$; and $4{,}208$ instances in
the $4$--$6$, $6$--$8$, and $8$--$10\,$m ranges respectively).
% TODO(figure): depth-range histogram + per-video (objects, length) scatter,
% analogous to prior group-depth datasets, once the plotting run is committed.

% VALIDATION (blocked on gate G1): a quantitative annotation-quality assessment
% -- mask IoU and per-instance depth-layer error against MANUAL ground truth on
% a hand-checked subset of ~150--300 frames -- is REQUIRED before any claim
% about annotation quality or the depth-layer MAE metric. Deliberately NOT
% written until that experiment exists; no placeholder numbers.

\section{Method}
\label{sec:method}
% Holistic depth initialization; instance depth-layer prediction; occlusion-
% aware refinement; single-sensor adaptations (metric SigLog, box-IoU pairing,
% visible-only supervision); temporal-consistency extension. BLOCKED: which
% components are claimed as contributions depends on G3/G8 + framing.

\subsection{Overview}
Our system follows the three-stage design of
InstanceDepth~\cite{liang2025instancedepth}, which we reproduce and adapt to the
single-sensor metric setting of Section~\ref{sec:problem} and extend with a
temporal-consistency module. Given a frame $I_t$: Stage~1
(Section~\ref{ssec:phase1}) predicts a dense metric depth map and multi-scale
depth features; Stage~2 (Section~\ref{ssec:phase2}) predicts a set of person
instances, each with a mask, a category, and a depth layer; and Stage~3
(Section~\ref{ssec:phase3}) identifies pairs of overlapping instances and
refines their depths so that occluder and occludee are ordered consistently.
Throughout, we retain the metric supervision afforded by the calibrated sensor
and adapt the occlusion machinery to modal, mutually-disjoint masks
(Section~\ref{ssec:adapt}); an optional streaming module then enforces temporal
consistency across frames (Section~\ref{ssec:temporal}). We describe the
reproduced components only to the level needed to make our adaptations precise,
citing InstanceDepth for full derivations, and state the components we introduce
explicitly as such.
% TODO(figure): system-overview figure (Fig. 1).

\subsection{Holistic Depth Initialization}
\label{ssec:phase1}
The first stage predicts a dense metric depth map together with the multi-scale
depth features reused by the later stages. Following
InstanceDepth~\cite{liang2025instancedepth}, the frame is encoded by a
pretrained backbone---in our experiments either a vanilla DINOv2
ViT-L/14~\cite{oquab2024dinov2} or the depth-fine-tuned encoder of Depth
Anything~V2~\cite{yang2024depthanythingv2} (compared in
Section~\ref{sec:ablation})---and a coarse-to-fine, patch-attention
\emph{depth-range decoder} produces per-level features $F_i$ at increasing
resolution.

Depth is then obtained by \emph{iterative bin refinement}. The metric range
$[0,\mathrm{MAX}_d]$ is divided into $r_d$ ordinal bins. Starting from an initial
estimate $D_0$, each level $i$ predicts a per-bin confidence
$C_i=\sigma(\Phi_d(F_i,D_i))$ and per-bin ordinal scores $S_i$, combines them
into a relative error score, and converts it into an additive metric
correction:
\begin{align}
\label{eq:refine}
R_i &= \sum_{b=0}^{r_d-1} C_{i,b}\,S_{i,b}, \\
E_i &= 2\,(R_i-1)\,\frac{\mathrm{MAX}_d}{r_d}, \qquad D_{i+1}=D_i+E_i .
\end{align}
Iterating across levels yields the dense metric depth $D$; we refer the reader
to InstanceDepth~\cite{liang2025instancedepth} for the full derivation of the
decoder and Eq.~\eqref{eq:refine}.

\smallskip\noindent\textbf{Adaptation: metric supervision.}
Because a single calibrated sensor provides true metric depth, we supervise $D$
with the scale-invariant logarithmic loss~\cite{eigen2014silog} \emph{without}
the scale-and-shift alignment used by relative-depth
methods~\cite{ranftl2022midas,yang2024depthanythingv2}. Writing the per-pixel
log-error as $\delta=\log D-\log D^{*}$ over the $n$ sensor-valid pixels
($V_t=1$),
\begin{equation}
\label{eq:silog}
\mathcal{L}_{\mathrm{depth}} =
\sqrt{\frac{1}{n}\sum \delta^2 \;-\; \lambda\Big(\frac{1}{n}\sum \delta\Big)^2},
\qquad \lambda=0.5 .
\end{equation}
The first term penalizes the magnitude of the log-error while the second removes
part of the influence of a global scale offset; setting $\lambda<1$ retains the
metric scale rather than discarding it, which is appropriate here and is
consistent with the metric depth-layer targets of Eq.~\eqref{eq:depthlayer}. The
loss is applied with deep supervision across the refinement levels.

\subsection{Instance Depth-Layer Prediction}
\label{ssec:phase2}
The second stage predicts, for each person, a segmentation mask, a category,
and an instance depth layer. Following InstanceDepth~\cite{liang2025instancedepth},
we use a Mask2Former~\cite{cheng2022mask2former} instance decoder with its own
Swin-L backbone~\cite{liu2021swin} (kept separate from the Stage-1 encoder so it
can be frozen during Stage-3 refinement): a fixed set of $N$ object queries is
decoded into per-query mask logits $\mathrm{Msk}_i$ and class scores
$\mathrm{Cls}_i$. To predict the instance depth layer $\mathrm{Dep}_i$ we attach
a lightweight head---a small MLP on each query embedding followed by a softplus
for positivity---mirroring the query-based class head. We use $N{=}100$ queries;
the COCO-pretrained default of $200$ is well above the number of people per
frame in our data (Table~\ref{tab:dataset}), and reducing it halves the decoder
and matching budget while re-using the pretrained weights unchanged apart from
the query embeddings.

Predictions are trained by bipartite matching against the ground-truth
instances~\cite{liang2025instancedepth,carion2020detr}. The optimal assignment
$\hat\sigma$ minimizes a cost that augments Mask2Former's mask and class terms
with a depth-layer term (Eq.~5--7 of~\cite{liang2025instancedepth}):
\begin{equation}
\label{eq:match}
\hat\sigma=\arg\min_{\sigma}\sum_i
\big(\lambda_m\,\mathcal{L}^{i}_{\mathrm{mask}}
    +\lambda_c\,\mathcal{L}^{i}_{\mathrm{cls}}
    +\lambda_d\,\mathcal{L}^{i}_{\mathrm{dep}}\big),
\end{equation}
where $\mathcal{L}^{i}_{\mathrm{mask}}$, $\mathcal{L}^{i}_{\mathrm{cls}}$, and
$\mathcal{L}^{i}_{\mathrm{dep}}$ measure the disagreement between query $i$ and
its assigned ground-truth instance $\sigma(i)$ in mask, class, and depth layer:
a point-sampled~\cite{kirillov2020pointrend} binary cross-entropy plus Dice, a
cross-entropy, and a smooth-$\ell_1$ regression, respectively. The same terms,
evaluated on the matched pairs---with unmatched queries supervised toward a
no-object class---form the training loss, and the ground-truth depth layer is
that of Eq.~\eqref{eq:depthlayer}.

\subsection{Occlusion-Aware Refinement}
\label{ssec:phase3}
The final stage corrects instance depths where people overlap. Following
InstanceDepth~\cite{liang2025instancedepth}, it (i) selects confident candidate
instances, (ii) forms occluder--occludee pairs, (iii) extracts paired evidence
by ROI alignment, and (iv) predicts a per-instance depth correction with a
relation-reasoning head.

\smallskip\noindent\textbf{Candidates and pairs.}
From the Stage-2 predictions we keep instances whose category and mask
confidences exceed $0.9$ and $0.8$~\cite{liang2025instancedepth}, and pair each
with its most-overlapping neighbour, ordering the pair so that the nearer
instance (the occluder) is first. \emph{Adaptation.} Because our ground-truth
masks are modal and mutually disjoint (Section~\ref{sec:problem}), two occluding
people have near-zero \emph{mask} overlap; we therefore detect overlap by
\emph{bounding-box} IoU (threshold $0.1$) rather than mask IoU.

\smallskip\noindent\textbf{Relation reasoning.}
For each pair, multi-scale depth features $F_{\mathrm{obj}}$ and geometric priors
$G_{\mathrm{obj}}$ (mask logits, normalized coordinates, and the base depth) are
extracted by ROI alignment~\cite{he2017maskrcnn} into a common grid, and a
relation head $\Phi_o$ predicts a bounded relative error that rescales the base
depth $D_{\mathrm{obj}}$ (Eq.~8--9 of~\cite{liang2025instancedepth}):
\begin{align}
\label{eq:eobj}
E_{\mathrm{obj}} &= \sigma\big(\Phi_o([F_{\mathrm{obj}}, G_{\mathrm{obj}}])\big),\\
\hat D_{\mathrm{obj}} &= (2E_{\mathrm{obj}}-1)\,D_{\mathrm{obj}} + D_{\mathrm{obj}}.
\end{align}
Because $\Phi_o$ receives both members of a pair, each instance's correction is
conditioned on its partner. As $E_{\mathrm{obj}}=\tfrac12$ gives
$\hat D_{\mathrm{obj}}=D_{\mathrm{obj}}$, the stage is initialized to the
identity.

\smallskip\noindent\textbf{Losses.}
The refined depths are supervised by a scale-invariant term and a
relative-ordering term (Eq.~10--12 of~\cite{liang2025instancedepth}):
\begin{align}
\mathcal{L}_{\mathrm{obj}} &= \mathrm{SILog}(\hat D_{\mathrm{obj}}, D^{*}_{\mathrm{obj}}),\\
\mathcal{L}_{\mathrm{dist}} &= \sum_{i\neq j}\big\lVert(\hat D_i-\hat D_j)^2-(D^{*}_i-D^{*}_j)^2\big\rVert,\\
\mathcal{L}_{\mathrm{ref}} &= \lambda_1\mathcal{L}_{\mathrm{obj}} + \lambda_2\mathcal{L}_{\mathrm{dist}}.
\end{align}
\emph{Adaptation.} A single sensor has no ground truth behind an occluder, so
$\mathcal{L}_{\mathrm{obj}}$ is evaluated only on the \emph{visible} pixels of
each matched instance; the relative term $\mathcal{L}_{\mathrm{dist}}$ pins the
depth gap between occluder and occludee to the ground-truth gap.
\emph{Non-degrading design.} We keep the Stage-1 depth branch \emph{frozen}
during refinement, so the dense map outside refined instances is left unchanged;
combined with the identity initialization, the stage can only alter the paired
instances and cannot degrade the base depth by construction.

\smallskip\noindent\textbf{Bounded pair-attention head (extension).}
We additionally study a variant of $\Phi_o$ that replaces its $1{\times}1$
convolutions with a small cross-attention encoder---so each instance's ROI
tokens attend to its partner's---and bounds the correction as
$\hat D_{\mathrm{obj}}=(1+c\tanh z)\,D_{\mathrm{obj}}$, limiting the
multiplicative rescaling to a fixed range $|c|$ rather than the
$(0,2)$ of Eq.~\eqref{eq:eobj}. It is a drop-in replacement with the same
interface; whether it improves over the original head is examined in
Section~\ref{sec:ablation}, not assumed here.

\subsection{Single-Sensor Adaptations}
\label{ssec:adapt}
The single-sensor, metric setting of Section~\ref{sec:problem} requires three
adaptations to the reproduced pipeline, which we collect here.
\emph{(i)~Metric supervision.} All depth losses use the scale-invariant log form
\emph{without} scale-and-shift alignment (Eq.~\eqref{eq:silog}), because the
calibrated sensor provides true metric depth that a scale-invariant objective
would discard. \emph{(ii)~Bounding-box occlusion pairing.} Since the ground-truth
and hence predicted instance masks are modal and mutually disjoint, occluding
people have near-zero mask overlap, so occlusion pairs are formed from
bounding-box rather than mask IoU. \emph{(iii)~Visible-only, non-degrading
refinement.} A single sensor observes only front surfaces, so the refinement
loss is masked to the visible pixels of each instance, and the Stage-1 depth
branch is frozen during refinement so that correcting occluded instances cannot
degrade the dense map elsewhere. Together these keep the method trainable, and
the refinement non-degrading, on single-sensor metric data.

\subsection{Temporal-Consistency Extension}
\label{ssec:temporal}
The pipeline so far is per-frame, so the depth of static structure can flicker
across a video (Section~\ref{sec:problem}). As an optional extension we insert a
lightweight streaming stabilizer between the Stage-1 decoder and its refinement
heads: a residual convolutional GRU~\cite{ballas2016convgru} that carries a
hidden state across frames on a downsampled feature grid, with a zero-initialized
output projection so that at initialization it is an exact identity and the
per-frame model is recovered. Its state is reset at sequence boundaries and
updated once per frame, giving constant memory and unbounded video length.

To reduce flicker without penalizing genuine motion, we train the stabilizer
with the temporal-gradient-matching (TGM) objective of Video Depth
Anything~\cite{chen2025vda}, which matches the prediction's frame-to-frame
log-depth change $\Delta d_t=\log D_t-\log D_{t-1}$ to the ground truth's
$\Delta d^{*}_t$ rather than warping frames by optical flow:
\begin{equation}
\label{eq:tgm}
\mathcal{L}_{\mathrm{tgm}} = \frac{1}{|\mathcal{V}|}\sum_{p\in\mathcal{V}}
\big|\,\Delta d_t - \Delta d^{*}_t\,\big|,
\end{equation}
over the pixels $\mathcal{V}$ valid in both frames. A prediction that follows the
ground-truth motion incurs no penalty; only flicker---change the ground truth
does not exhibit---is penalized. The stabilizer is trained on short clips with
the spatial branch frozen, under the per-frame loss (Eq.~\eqref{eq:silog}) plus a
weighted $\mathcal{L}_{\mathrm{tgm}}$.

\section{Experimental Setup}
\label{sec:setup}
% Splits, resolutions, training schedules, hardware, and exact configuration
% profiles. BLOCKED: must match the final resolved code (G0) and state the
% enhanced-vs-faithful profile per run (G5).
\emph{[Written in the Experimental-Setup phase.]}

\section{Evaluation Metrics}
\label{sec:metrics}
% Depth metrics (AbsRel, RMS, RMS-log, log10, sigma_1..3); instance metrics
% (precision/recall/F1, mIoU, depth-layer MAE); temporal metrics (TAE, flicker
% ratio). BLOCKED: the occlusion-slice metric definition must be stated
% precisely (G6), and the eval protocol unified across phases (G7).
\emph{[Written in the Metrics phase.]}

\section{Results}
\label{sec:results}
% Main quantitative tables + qualitative figures. BLOCKED: G0 (numbers must
% come from code that matches the manuscript), G2 (significance), G4 (a fair
% external baseline). No numbers are written until these hold.
\emph{[Written in the Results phase.]}

\section{Ablation Studies}
\label{sec:ablation}
% Encoder (vanilla DINOv2 vs Depth-Anything-V2); regression loss (SigLog vs
% berHu); occlusion head (vanilla vs proposed); temporal loss on/off;
% profile ablations. BLOCKED: G3 (head attribution), G8 (temporal isolation).
\emph{[Written in the Ablation phase.]}

\section{Discussion}
\label{sec:discussion}
\emph{[Written in the Discussion phase.]}

\section{Limitations}
\label{sec:limitations}
% Single-sensor visible-only supervision; annotation reliance on model output;
% modest, dataset-specific gains; no amodal ground truth.
\emph{[Written in the Limitations phase.]}

\section{Conclusion}
\label{sec:conclusion}
\emph{[Written in the Conclusion phase.]}

%------------------------------------------------------------------------------
%  REFERENCES
%------------------------------------------------------------------------------
\bibliographystyle{ACM-Reference-Format}
\bibliography{references}

\end{document}
